\chapter{The Complete Solution to Ax = b}

So far you have learned to ask two separate questions about a matrix $\mat{A}$. The
column space told you which right-hand sides $\vec{b}$ are reachable: $\mat{A}\vec{x}=\vec{b}$
has \emph{some} solution exactly when $\vec{b}$ lies in $\col(\mat{A})$. The nullspace told
you about the slack: $\mat{A}\vec{x}=\vec{0}$ has nontrivial solutions exactly when the
columns are dependent. In this chapter we glue those two ideas together and answer the
question that systems of equations were invented to answer: \emph{given a matrix $\mat{A}$
and a vector $\vec{b}$, describe every $\vec{x}$ with $\mat{A}\vec{x}=\vec{b}$ -- all of
them, no more and no fewer.}

That full description is called the \emph{complete solution}, and it has a shape you can
draw. It is always

\[
  \vec{x} \;=\; \underbrace{\vec{x}_p}_{\text{one particular solution}}
        \;+\; \underbrace{\vec{x}_n}_{\text{anything in the nullspace}}.
\]

One special solution, plus the entire nullspace, slid over to sit on top of it. The rest of
the chapter unpacks that one line: why it is true, how to compute each piece, and what the
geometry looks like in every case (no solution, one solution, infinitely many).

\section{One picture to hold onto}

Before any algebra, here is the mental image. The map $\vec{x}\mapsto \mat{A}\vec{x}$ takes
the input space and folds it down onto the column space. Solving $\mat{A}\vec{x}=\vec{b}$
means: find every input that lands exactly on the target point $\vec{b}$.

\begin{intuitionbox}
Think of $\mat{A}$ as a machine that flattens. If two different inputs $\vec{x}_1$ and
$\vec{x}_2$ both produce $\vec{b}$, then their difference satisfies
\[
  \mat{A}(\vec{x}_1-\vec{x}_2)=\mat{A}\vec{x}_1-\mat{A}\vec{x}_2=\vec{b}-\vec{b}=\vec{0}.
\]
So the difference lives in the nullspace. Two solutions can differ \emph{only} by a
nullspace vector. Turn that around: take one solution $\vec{x}_p$ and add any nullspace
vector $\vec{x}_n$; then $\mat{A}(\vec{x}_p+\vec{x}_n)=\vec{b}+\vec{0}=\vec{b}$, so you get
another solution. \textbf{The solution set is the nullspace, picked up and translated to
pass through $\vec{x}_p$.}
\end{intuitionbox}

\begin{figure}[H]
\centering
\begin{tikzpicture}[scale=1.0,>=Stealth]
  % nullspace line through origin
  \draw[spacegray,thick] (-2.2,-2.2) -- (2.2,2.2);
  \node[spacegray, above left] at (1.8,1.9) {$\nullsp(\mat{A})$};
  \fill (0,0) circle (1.6pt) node[below left]{\small $\vec 0$};
  % solution set: parallel line through x_p
  \draw[vecblue,very thick] (-1.4,-2.2) -- (3.0,2.2);
  \node[vecblue, above right] at (2.8,2.2) {solution set of $\mat{A}\vec x=\vec b$};
  % particular solution
  \draw[->,vecred,very thick] (0,0) -- (1.3,0.5);
  \fill[vecred] (1.3,0.5) circle (1.8pt) node[below right]{\small $\vec x_p$};
  % a nullspace vector added
  \draw[->,spacegray,thick,dashed] (1.3,0.5) -- (2.3,1.5);
  \fill (2.3,1.5) circle (1.6pt) node[right]{\small $\vec x_p+\vec x_n$};
  \node[spacegray] at (2.15,0.85) {\small $\vec x_n$};
\end{tikzpicture}
\caption{The complete solution set (blue) is the nullspace (gray, through the origin)
slid over so it passes through one particular solution $\vec x_p$. Adding any nullspace
vector $\vec x_n$ to $\vec x_p$ keeps you on the blue line.}
\label{ch13:translate}
\end{figure}

A set like the blue line in Figure~\ref{ch13:translate} -- a subspace shifted off the
origin -- is called an \emph{affine subspace} (or a \emph{coset}). It is parallel to the
nullspace but generally does \emph{not} pass through $\vec{0}$. The only time it passes
through the origin is when $\vec{b}=\vec{0}$, because then $\vec{x}_p=\vec{0}$ works and the
solution set \emph{is} the nullspace.

\begin{focusbox}
Two questions, two subspaces, one formula.
\begin{itemize}
  \item \textbf{Existence} (does a solution exist?) is a question about the \emph{column
    space}: solvable $\iff \vec{b}\in\col(\mat{A})$.
  \item \textbf{Uniqueness} (is it the only one?) is a question about the
    \emph{nullspace}: unique $\iff \nullsp(\mat{A})=\{\vec{0}\}$.
  \item When both succeed for a given $\vec{b}$, the answer set is
    $\vec{x}_p+\nullsp(\mat{A})$.
\end{itemize}
\end{focusbox}

\section{Why particular plus nullspace is everything}

Now turn the picture into a theorem. The argument is short, and it is the same one that
reappears for differential equations, recurrences, and every other linear problem you will
meet.

\begin{theorem}[Structure of the solution set]
\label{ch13:structure}
Suppose $\mat{A}\vec{x}_p=\vec{b}$ for some particular vector $\vec{x}_p$ (that is,
$\vec{b}$ is reachable). Then the complete set of solutions to $\mat{A}\vec{x}=\vec{b}$ is
\[
  \set{\vec{x}_p+\vec{x}_n \;:\; \vec{x}_n\in\nullsp(\mat{A})}.
\]
If instead no particular solution exists, the solution set is empty.
\end{theorem}

\begin{proof}
We prove the two halves of set equality.

\emph{Every $\vec{x}_p+\vec{x}_n$ is a solution.} If $\mat{A}\vec{x}_n=\vec{0}$, then by
linearity
\[
  \mat{A}(\vec{x}_p+\vec{x}_n)=\mat{A}\vec{x}_p+\mat{A}\vec{x}_n=\vec{b}+\vec{0}=\vec{b}.
\]

\emph{Every solution has this form.} Suppose $\mat{A}\vec{x}=\vec{b}$. Set
$\vec{x}_n:=\vec{x}-\vec{x}_p$. Then
\[
  \mat{A}\vec{x}_n=\mat{A}\vec{x}-\mat{A}\vec{x}_p=\vec{b}-\vec{b}=\vec{0},
\]
so $\vec{x}_n\in\nullsp(\mat{A})$, and $\vec{x}=\vec{x}_p+\vec{x}_n$ is of the claimed
form. The two inclusions together give equality.
\end{proof}

Notice what the proof did \emph{not} need: it never used the size of $\mat{A}$, never
assumed a square matrix, never assumed invertibility. It used only linearity. That is why
the formula is universal.

\begin{trapbox}
The complete solution is $\vec{x}_p+\nullsp(\mat{A})$, \textbf{not} $\nullsp(\mat{A})$ and
\textbf{not} ``the span of $\vec{x}_p$ and the nullspace.'' You do not scale $\vec{x}_p$.
You add it once, as a fixed offset, and let only the nullspace part vary. Writing
$c\vec{x}_p+(\text{nullspace})$ is wrong: $2\vec{x}_p$ usually solves
$\mat{A}\vec{x}=2\vec{b}$, not $\mat{A}\vec{x}=\vec{b}$.
\end{trapbox}

\begin{strangbox}
Strang states this as ``\emph{complete solution} $=$ \emph{particular} $+$
\emph{special}.'' His ``particular solution'' is our $\vec{x}_p$, gotten by setting all the
free variables to zero. His ``special solutions'' $\vec{s}_1,\vec{s}_2,\dots$ are a basis
for the nullspace, gotten by setting one free variable to $1$ and the rest to $0$. So
$\vec{x}_n=c_1\vec{s}_1+\cdots+c_k\vec{s}_k$, and the complete solution is
$\vec{x}_p+c_1\vec{s}_1+\cdots+c_k\vec{s}_k$. Reading this chapter \emph{before} Strang,
you now know \emph{why} that recipe produces everything; reading it \emph{after}, you can
see it is Theorem~\ref{ch13:structure} with the nullspace written in a basis.
\end{strangbox}

\section{Existence: is \texorpdfstring{$\vec b$}{b} in the column space?}

A solution exists if and only if $\vec{b}=\mat{A}\vec{x}$ for some $\vec{x}$ -- and
$\mat{A}\vec{x}$ is, by definition, a linear combination of the columns of $\mat{A}$. So
the reachable right-hand sides are exactly the column space. There is nothing more to it
conceptually; the work is in \emph{detecting} membership, which elimination does for you.

When you row-reduce the augmented matrix $[\,\mat{A}\mid\vec{b}\,]$, watch the rows that
become all zeros on the $\mat{A}$ side. If such a zero row has a nonzero entry on the
$\vec{b}$ side, you have produced the equation $0=(\text{nonzero})$, which is a
contradiction: $\vec{b}$ was \emph{not} in the column space, and there is no solution. If
every all-zero row of $\mat{A}$ pairs with a zero in $\vec{b}$, the system is consistent.

\begin{figure}[H]
\centering
\begin{tikzpicture}[scale=1.0,>=Stealth]
  % a plane = column space in 3D-ish view
  \begin{scope}[shift={(0,0)}]
    \fill[targetgreen!12] (-1.6,-0.9) -- (2.4,-0.4) -- (2.0,1.6) -- (-2.0,1.1) -- cycle;
    \draw[targetgreen!70] (-1.6,-0.9) -- (2.4,-0.4) -- (2.0,1.6) -- (-2.0,1.1) -- cycle;
    \node[targetgreen!60!black] at (1.5,1.35) {\small $\col(\mat A)$};
    \fill (0,0) circle (1.4pt) node[below left]{\small $\vec0$};
    % b in the plane
    \draw[->,vecblue,very thick] (0,0) -- (1.1,0.35);
    \node[vecblue] at (1.25,0.05) {\small $\vec b_{\text{in}}$};
    % b off the plane
    \draw[->,vecred,very thick] (0,0) -- (0.4,1.9);
    \node[vecred] at (0.55,2.05) {\small $\vec b_{\text{off}}$};
    \draw[vecred,dashed] (0.4,1.9) -- (0.0,0.55);
  \end{scope}
\end{tikzpicture}
\caption{Existence is a hit-or-miss question. If $\vec b$ lies in the plane
$\col(\mat A)$ (blue), the system is solvable. If $\vec b$ sticks up out of the plane
(red), no input can reach it and $\mat A\vec x=\vec b$ has no solution.}
\label{ch13:existence}
\end{figure}

\begin{example}[Spotting an inconsistent system]
\label{ch13:incon}
Take
\[
  \mat{A}=\begin{bmatrix}1&2\\2&4\\3&6\end{bmatrix},\qquad
  \vec{b}=\colvec{1\\3\\3}.
\]
Every column of $\mat{A}$ is a multiple of $\colvec{1\\2\\3}$, so the column space is the
single line through $\colvec{1\\2\\3}$. Is $\vec{b}$ on that line? It would need
$\vec{b}=t\colvec{1\\2\\3}$, forcing $t=1$ from the first entry but $t=\tfrac32$ from the
second. No single $t$ works, so $\vec{b}\notin\col(\mat{A})$ and there is \emph{no}
solution. Elimination confirms it: subtract $2\times$row~1 from row~2 and $3\times$row~1
from row~3 in $[\,\mat{A}\mid\vec{b}\,]$:
\[
  \left[\begin{array}{cc|c}1&2&1\\2&4&3\\3&6&3\end{array}\right]
  \;\longrightarrow\;
  \left[\begin{array}{cc|c}1&2&1\\0&0&1\\0&0&0\end{array}\right].
\]
Row~2 now reads $0x_1+0x_2=1$, i.e.\ $0=1$. Contradiction; inconsistent.
\end{example}

\begin{example}[Same $\mat A$, a reachable $\vec b$]
\label{ch13:con}
Keep the same $\mat{A}$ but use $\vec{b}=\colvec{2\\4\\6}=2\colvec{1\\2\\3}$, which
\emph{is} on the column line. Reducing,
\[
  \left[\begin{array}{cc|c}1&2&2\\2&4&4\\3&6&6\end{array}\right]
  \;\longrightarrow\;
  \left[\begin{array}{cc|c}1&2&2\\0&0&0\\0&0&0\end{array}\right].
\]
Now the zero rows of $\mat{A}$ pair with zeros in $\vec{b}$: consistent. The one surviving
equation is $x_1+2x_2=2$, which has infinitely many solutions -- a whole line of them. We
will write that line out in the next section.
\end{example}

\begin{walkbox}
On your walk, fix the $3\times2$ matrix $\mat{A}$ from the examples in your head -- its
column space is one line in $\R^3$. Now imagine sweeping $\vec{b}$ all over $\R^3$. For how
many $\vec{b}$ is the system solvable: none, a line's worth, a plane's worth, or all of
$\R^3$? And of the solvable ones, how many have a \emph{unique} solution? (Hint: the
nullspace of this $\mat A$ is not just the origin.)
\end{walkbox}

\section{Uniqueness: is the nullspace trivial?}

Existence was about $\vec{b}$ and the columns reaching far enough. Uniqueness is about the
columns being non-redundant. By Theorem~\ref{ch13:structure}, once a solution exists, the
number of solutions equals the number of nullspace vectors:

\[
  \nullsp(\mat{A})=\{\vec{0}\}\;\Longrightarrow\;\text{at most one solution};\qquad
  \nullsp(\mat{A})\neq\{\vec{0}\}\;\Longrightarrow\;\text{either none or infinitely many}.
\]

There is never a system with exactly two or exactly seventeen solutions. The solution set
is an affine copy of a subspace, and a nonzero subspace already contains infinitely many
vectors (you can scale any nonzero element). So the count is always $0$, $1$, or $\infty$.

\begin{aestheticbox}
This trichotomy -- zero, one, or infinitely many -- is one of the things linearity buys
you. Nonlinear equations are wild: $x^2=c$ can have $0$, $1$, or $2$ roots; $\sin
x=c$ can have infinitely many \emph{isolated} roots. Linear systems are tame because
straightness is preserved: if two solutions exist, the entire line through them solves the
system too. The geometry of the answer set is always flat. Hold onto this; it is the reason
linear models are the workhorses of science.
\end{aestheticbox}

\begin{trapbox}
``Square matrix'' does not mean ``unique solution.'' A $3\times3$ matrix can have a
nontrivial nullspace (when its columns are dependent), giving zero or infinitely many
solutions. Uniqueness is controlled by the nullspace, which is controlled by
\emph{independence of the columns}, not by the matrix being square. We will see all four
combinations below.
\end{trapbox}

\section{Rank decides the shape}

Everything -- existence, uniqueness, and the dimension of the solution set -- is governed
by one number and the two dimensions of $\mat{A}$. Let $\mat{A}$ be $m\times n$ (that is,
$m$ equations in $n$ unknowns) with rank $r$. Recall the two facts you have already proved
about $r$:

\begin{itemize}
  \item $r=\dim\col(\mat{A})$: the number of pivots, the number of independent columns,
    the dimension of the reachable target space.
  \item $\dim\nullsp(\mat{A})=n-r$: the number of free (non-pivot) columns. This is the
    \emph{rank--nullity} count, and it is exactly the number of free variables in the
    solution.
\end{itemize}

Combining these with Theorem~\ref{ch13:structure} gives the master statement of the
chapter.

\begin{formalbox}
\textbf{The shape of the solution set.} For $\mat{A}\in\R^{m\times n}$ of rank $r$ and a
\emph{consistent} right-hand side $\vec{b}$, the solution set is an affine subspace of
dimension $n-r$:
\[
  \vec{x}_p+\nullsp(\mat{A}),\qquad \dim=n-r.
\]
In particular:
\begin{itemize}
  \item solvable for \emph{every} $\vec{b}\in\R^m$ $\iff$ $r=m$ (the columns span $\R^m$;
    \emph{full row rank});
  \item solution unique whenever it exists $\iff$ $r=n$ (the columns are independent;
    \emph{full column rank}).
\end{itemize}
\end{formalbox}

Read the two bullets as ``existence for all $\vec b$ is about rows'' and ``uniqueness is
about columns.'' The number $r$ can be at most $\min(m,n)$, so a matrix can be full row
rank, full column rank, both (then it is square and invertible), or neither.

\begin{figure}[H]
\centering
\begin{tikzpicture}[scale=0.92,>=Stealth]
  % Column headers
  \node[draw=spacegray!60, fill=gray!10, rounded corners=2pt, font=\footnotesize\scshape, text=black!85, minimum width=5.2cm, align=center]
    at (0, 2.8) {Full Row Rank ($r = m$)\\[1pt]\scriptsize (Existence for all $\vec{b}$)};
  \node[draw=spacegray!60, fill=gray!10, rounded corners=2pt, font=\footnotesize\scshape, text=black!85, minimum width=5.2cm, align=center]
    at (5.8, 2.8) {Not Full Row Rank ($r < m$)\\[1pt]\scriptsize (Zero rows in $\mat{R}$; solvability restricted)};

  % Row headers
  \node[draw=spacegray!60, fill=gray!10, rounded corners=2pt, font=\footnotesize\scshape, text=black!85, minimum height=2.2cm, align=center, text width=2.3cm]
    at (-4.3, 1.1) {Full Column Rank\\[2pt]($r = n$)\\[2pt]\scriptsize Uniqueness};
  \node[draw=spacegray!60, fill=gray!10, rounded corners=2pt, font=\footnotesize\scshape, text=black!85, minimum height=2.2cm, align=center, text width=2.3cm]
    at (-4.3, -1.5) {Not Full Column Rank\\[2pt]($r < n$)\\[2pt]\scriptsize Free variables; non-unique};

  % Cell (1,1): Square & Invertible
  \node[draw=targetgreen!80, fill=targetgreen!12, rounded corners=3pt, minimum width=5.2cm, minimum height=2.2cm, align=center, font=\footnotesize]
    at (0, 1.1) {
      \textbf{\color{targetgreen!85!black}$r = m = n$ (Square Invertible)}\\[3pt]
      \scriptsize Nullspace $= \{\vec{0}\}$, Columns span $\R^m$\\[2pt]
      \textbf{Exactly $1$ unique solution} for every $\vec{b}$:\\[2pt]
      $\vec{x} = \mat{A}^{-1}\vec{b}$
    };

  % Cell (1,2): Tall / Full Column Rank
  \node[draw=sourceblue!80, fill=sourceblue!12, rounded corners=3pt, minimum width=5.2cm, minimum height=2.2cm, align=center, font=\footnotesize]
    at (5.8, 1.1) {
      \textbf{\color{sourceblue!85!black}$r = n < m$ (Tall, Independent)}\\[3pt]
      \scriptsize Nullspace $= \{\vec{0}\}$, Column space is proper\\[2pt]
      \textbf{$0$ or $1$ solution}:\\[2pt]
      $1$ if $\vec{b} \in \col(\mat{A})$, $0$ if $\vec{b} \notin \col(\mat{A})$
    };

  % Cell (2,1): Wide / Full Row Rank
  \node[draw=quotientorange!80, fill=quotientorange!14, rounded corners=3pt, minimum width=5.2cm, minimum height=2.2cm, align=center, font=\footnotesize]
    at (0, -1.5) {
      \textbf{\color{quotientorange!90!black}$r = m < n$ (Wide, Spanning)}\\[3pt]
      \scriptsize $n - m$ free variables; columns span $\R^m$\\[2pt]
      \textbf{Infinitely many solutions} ($\infty$) for every $\vec{b}$:\\[2pt]
      affine flat $\vec{x}_p + \nullsp(\mat{A})$
    };

  % Cell (2,2): General Case
  \node[draw=trapred!80, fill=trapred!12, rounded corners=3pt, minimum width=5.2cm, minimum height=2.2cm, align=center, font=\footnotesize]
    at (5.8, -1.5) {
      \textbf{\color{trapred!85!black}$r < m$ and $r < n$ (General Case)}\\[3pt]
      \scriptsize Zero rows in $\mat{R}$ \emph{and} free variables\\[2pt]
      \textbf{$0$ or $\infty$ solutions}:\\[2pt]
      $0$ if inconsistent, $\infty$ if $\vec{b} \in \col(\mat{A})$
    };
\end{tikzpicture}
\caption{The four cases of the linear system $\mat{A}\vec{x}=\vec{b}$, categorized by row rank
$r = m$ (which guarantees that columns span $\R^m$, ensuring existence) and column rank $r = n$
(which guarantees that $\nullsp(\mat{A}) = \{\vec{0}\}$, ensuring uniqueness).}
\label{ch13:fourcases}
\end{figure}

\subsection{The four cases, one at a time}

\textbf{Case 1: $r=m=n$ (square and invertible).} Full row rank gives existence for every
$\vec{b}$; full column rank gives uniqueness. Together: \emph{exactly one} solution for
every $\vec{b}$, namely $\vec{x}=\mat{A}\inv\vec{b}$. The nullspace is $\{\vec 0\}$, so
$\vec{x}_p$ is the whole story. These are the comfortable square invertible systems.

\textbf{Case 2: $r=n<m$ (tall, full column rank).} More equations than unknowns. The
columns are independent, so the nullspace is trivial and any solution is unique -- but the
columns may not reach every $\vec b$, so existence can fail. Outcome: \emph{$0$ or $1$}
solution. Either $\vec{b}$ lies in the $n$-dimensional column space (one solution) or it
does not (none). This is the setting of \emph{least squares}: when no exact solution
exists, you will later settle for the closest reachable point.

\textbf{Case 3: $r=m<n$ (wide, full row rank).} More unknowns than equations. The columns
span all of $\R^m$, so existence holds for \emph{every} $\vec{b}$ -- but there are
$n-m>0$ free variables, so the nullspace is nonzero and solutions are never unique.
Outcome: \emph{$\infty$} solutions for every $\vec{b}$. These are the
\emph{underdetermined} systems, where the minimum-norm solution earns its keep.

\textbf{Case 4: $r<m$ and $r<n$ (the general rectangular case).} Neither full. Existence
can fail \emph{and} uniqueness fails whenever existence holds. Outcome: \emph{$0$ or
$\infty$}. Almost every real-world data matrix is here.

\begin{labbox}
Build a tiny example of each case with your own hands and confirm the count.
\begin{enumerate}
  \item Case~1: $\mat{A}=\begin{bmatrix}2&0\\0&3\end{bmatrix}$, $\vec b=\colvec{4\\3}$.
    Solve; you should get the single point $\colvec{2\\1}$.
  \item Case~2: $\mat{A}=\begin{bmatrix}1&0\\0&1\\1&1\end{bmatrix}$. Try
    $\vec{b}=\colvec{1\\2\\3}$ (consistent) and $\vec{b}=\colvec{1\\2\\4}$ (inconsistent).
    What is special about the first $\vec b$?
  \item Case~3: $\mat{A}=\begin{bmatrix}1&1&1\end{bmatrix}$, $\vec{b}=[2]$. Write the whole
    plane of solutions in the form $\vec{x}_p+\nullsp(\mat A)$.
  \item Case~4: $\mat{A}=\begin{bmatrix}1&2\\2&4\end{bmatrix}$. Find a $\vec b$ with no
    solution and a $\vec b$ with infinitely many.
\end{enumerate}
\end{labbox}

\section{A full worked example: building the complete solution}

Let us run the whole machine once, slowly, on a matrix that is neither square nor full
rank, so you see every moving part. Take

\[
  \mat{A}=\begin{bmatrix}1&3&0&2\\0&0&1&4\\1&3&1&6\end{bmatrix},
  \qquad
  \vec{b}=\colvec{1\\6\\7}.
\]

This is $3\times4$: three equations, four unknowns. We expect free variables.

\paragraph{Step 1 -- eliminate.} Work on the augmented matrix. Subtract row~1 from row~3:
\[
  \left[\begin{array}{cccc|c}
    1&3&0&2&1\\ 0&0&1&4&6\\ 1&3&1&6&7
  \end{array}\right]
  \xrightarrow{\,R_3-R_1\,}
  \left[\begin{array}{cccc|c}
    1&3&0&2&1\\ 0&0&1&4&6\\ 0&0&1&4&6
  \end{array}\right]
  \xrightarrow{\,R_3-R_2\,}
  \left[\begin{array}{cccc|c}
    1&3&0&2&1\\ 0&0&1&4&6\\ 0&0&0&0&0
  \end{array}\right].
\]
The last row is all zeros, including the augmented entry: \textbf{consistent}. The pivots
sit in columns $1$ and $3$, so the rank is $r=2$. Columns $2$ and $4$ have no pivot --
those are the \emph{free} columns, and $x_2,x_4$ are the \emph{free variables}. The
nullspace will have dimension $n-r=4-2=2$.

\paragraph{Step 2 -- a particular solution (free variables $=0$).} Set $x_2=0$ and
$x_4=0$. The reduced equations read
\[
  x_1+3x_2+2x_4 = 1 \;\Rightarrow\; x_1=1,\qquad
  x_3+4x_4 = 6 \;\Rightarrow\; x_3=6.
\]
So
\[
  \vec{x}_p=\colvec{1\\0\\6\\0}.
\]
Sanity check: $\mat{A}\vec{x}_p$ should be $\vec{b}$. First entry
$1\cdot1+3\cdot0+0\cdot6+2\cdot0=1$; second $0+0+1\cdot6+0=6$; third
$1+0+6+0=7$. Indeed $\colvec{1\\6\\7}=\vec b$. Good.

\paragraph{Step 3 -- the special solutions (nullspace basis).} Now solve
$\mat{A}\vec{x}=\vec{0}$. Same pivots; the homogeneous equations are
\[
  x_1=-3x_2-2x_4,\qquad x_3=-4x_4.
\]
Turn each free variable on by itself.

Set $x_2=1,\ x_4=0$: then $x_1=-3,\ x_3=0$, giving
$\vec{s}_1=\colvec{-3\\1\\0\\0}$.

Set $x_2=0,\ x_4=1$: then $x_1=-2,\ x_3=-4$, giving
$\vec{s}_2=\colvec{-2\\0\\-4\\1}$.

These two special solutions are a basis for the nullspace, so
$\nullsp(\mat{A})=\spn\{\vec{s}_1,\vec{s}_2\}$ -- a $2$-dimensional plane through the
origin in $\R^4$, exactly as the count $n-r=2$ predicted.

\paragraph{Step 4 -- assemble the complete solution.} By
Theorem~\ref{ch13:structure},
\[
  \boxed{\;
  \vec{x}=\colvec{1\\0\\6\\0}
     + c_1\colvec{-3\\1\\0\\0}
     + c_2\colvec{-2\\0\\-4\\1},
  \qquad c_1,c_2\in\R.\;}
\]
That is \emph{every} solution: a particular point in $\R^4$ plus a $2$-dimensional plane of
freedom. Pick any $c_1,c_2$ and you get a valid $\vec{x}$; conversely every valid $\vec{x}$
arises from some choice. There are infinitely many solutions, organized as a flat
$2$-dimensional sheet.

\begin{remark}
Notice the bookkeeping that made this painless: the \emph{same} elimination served both
jobs. The pivots told us the rank and which variables are free; setting free variables to
$0$ produced $\vec{x}_p$; turning them on one at a time produced the nullspace basis. You
never reduce twice.
\end{remark}

\begin{trapbox}
A frequent slip: forgetting to put the free variables \emph{into} the special solutions.
When you set $x_2=1$, the special solution's second entry must be $1$, not $0$ -- the free
variable is part of the vector. Each special solution has a single $1$ in its own free slot
and $0$ in the other free slots; the pivot slots are then forced. If your special solutions
do not have those isolated $1$s, recheck.
\end{trapbox}

\section{Reading the four cases off your worked example}

The example above had $m=3$, $n=4$, $r=2$. Since $r<m$ ($2<3$) it is not full row rank, so
some right-hand sides have \emph{no} solution. Since $r<n$ ($2<4$) it is not full column
rank, so when a solution exists it is never unique. That is Case~4. Our chosen $\vec{b}$
happened to be consistent, so we landed in the ``$\infty$'' branch. Had we used, say,
$\vec{b}=\colvec{1\\6\\99}$, the last augmented row would have read $0=92$ and we would
have been in the ``$0$'' branch.

\begin{example}[Flipping a case by changing the shape]
\label{ch13:flip}
Delete the last column of the worked $\mat{A}$ to get the $3\times3$ matrix
$\mat{A}'=\begin{bmatrix}1&3&0\\0&0&1\\1&3&1\end{bmatrix}$. Column~2 is still a non-pivot
(it equals $3\times$ column~1), so $r=2$ \emph{again}, and now $r<n=3$ and $r<m=3$: still
Case~4, still $0$ or $\infty$. Squareness did not buy uniqueness, because the columns are
still dependent. To reach Case~1 you must make the columns independent, e.g.\ replace
column~2 by $\colvec{0\\1\\0}$, which gives a pivot in every column.
\end{example}

\section{The two ``nice'' rectangular extremes}

The full-row-rank and full-column-rank cases are worth a moment alone, because they are the
backbone of regression and of underdetermined inverse problems. Figure~\ref{ch13:extremes}
contrasts their geometry: a tall full-column-rank matrix can \emph{miss} $\vec b$ but never
offers a choice, while a wide full-row-rank matrix always \emph{reaches} $\vec b$ but always
offers a whole line (or more) of choices.

\begin{figure}[H]
\centering
\begin{tikzpicture}[scale=1.0,>=Stealth]
  % LEFT: full column rank -- point off a plane in R^3
  \begin{scope}[shift={(0,0)}]
    \node[font=\small\bfseries] at (0,2.5) {full column rank ($r=n$)};
    \fill[targetgreen!12] (-1.7,-0.9) -- (1.9,-0.5) -- (1.6,1.4) -- (-2.0,1.0) -- cycle;
    \draw[targetgreen!70] (-1.7,-0.9) -- (1.9,-0.5) -- (1.6,1.4) -- (-2.0,1.0) -- cycle;
    \node[targetgreen!60!black] at (1.0,1.2) {\small $\col(\mat A)$};
    \fill (0,0) circle (1.2pt);
    \draw[->,vecblue,very thick] (0,0) -- (0.9,0.2);
    \node[vecblue] at (1.15,0.0) {\small reachable};
    \draw[->,vecred,very thick] (0,0) -- (-0.3,1.8);
    \node[vecred] at (-0.3,2.05) {\small $\vec b$ off plane};
    \draw[vecred,dashed] (-0.3,1.8) -- (-0.6,0.45);
    \node[align=center,font=\small] at (0,-1.6) {$0$ or $1$ solution\\(never a choice)};
  \end{scope}
  % RIGHT: full row rank -- line of solutions
  \begin{scope}[shift={(6.4,0)}]
    \node[font=\small\bfseries] at (0,2.5) {full row rank ($r=m$)};
    \draw[spacegray,thick] (-1.8,-1.4) -- (1.8,1.4);
    \node[spacegray] at (1.7,1.65) {\small $\nullsp(\mat A)$};
    \fill (0,0) circle (1.2pt) node[below left]{\small $\vec0$};
    \draw[vecblue,very thick] (-1.0,-1.7) -- (2.6,1.1);
    \node[vecblue] at (2.5,1.65) {\small solution line};
    \draw[->,vecred,very thick] (0,0) -- (0.8,-0.3);
    \fill[vecred] (0.8,-0.3) circle (1.6pt) node[below]{\small $\vec x_p$};
    \node[align=center,font=\small] at (0.4,-1.85) {always $\infty$ solutions\\(for every $\vec b$)};
  \end{scope}
\end{tikzpicture}
\caption{The two rectangular extremes. \emph{Left:} full column rank --- the column space
may not contain $\vec b$ (no solution), but when it does the answer is a single point, since
the nullspace is just $\{\vec 0\}$. \emph{Right:} full row rank --- every $\vec b$ is
reachable, but the nonzero nullspace slides $\vec x_p$ into a whole flat of solutions.}
\label{ch13:extremes}
\end{figure}

\subsection{Full column rank: at most one solution}

If $\rank(\mat{A})=n$ (every column is a pivot column), there are no free variables, so the
nullspace is just $\{\vec{0}\}$. The complete solution collapses to a single point
$\vec{x}_p$ whenever $\vec{b}$ is reachable. The tall, thin matrix
\[
  \mat{A}=\begin{bmatrix}1&0\\1&1\\1&2\end{bmatrix}
\]
has two independent columns ($r=2=n$). For $\vec{b}=\colvec{1\\2\\3}$ you can check
$\vec{x}=\colvec{1\\1}$ works and is the \emph{only} solution. For
$\vec{b}=\colvec{1\\2\\4}$, elimination produces a contradiction and there is no exact
solution at all -- the point sticks up out of the $2$-dimensional column plane in $\R^3$.

\subsection{Full row rank: always at least one solution}

If $\rank(\mat{A})=m$ (every row has a pivot), the columns span all of $\R^m$, so every
$\vec{b}$ is reachable and the system is \emph{always} consistent. The wide matrix
\[
  \mat{A}=\begin{bmatrix}1&1&1\\0&1&2\end{bmatrix}
\]
has $r=2=m$. Both rows pivot, so for any $\vec b\in\R^2$ a solution exists; and since
$n-r=3-2=1$, the solution set is always a line. For $\vec b=\colvec{3\\3}$ one particular
solution is found by setting the lone free variable $x_3=0$: back-substitution gives
$x_2=3,\ x_1=0$, so $\vec{x}_p=\colvec{0\\3\\0}$. The nullspace special solution (set
$x_3=1$) is $\vec{s}=\colvec{1\\-2\\1}$, so the complete solution is the line
$\colvec{0\\3\\0}+c\colvec{1\\-2\\1}$.

\begin{connectionbox}
The split ``$\vec{x}_p+$ homogeneous solutions'' is not special to matrices -- it is the
signature of \emph{any} linear operator $L$ with $L\vec{x}=\vec{b}$. In a differential
equations course you meet exactly this when solving $y''+y=\cos t$: the general solution is
one \emph{particular} solution (a response forced by $\cos t$) plus the general solution of
the \emph{homogeneous} equation $y''+y=0$ (the nullspace of the operator
$L=\tfrac{d^2}{dt^2}+1$, namely $c_1\cos t+c_2\sin t$). Same theorem, same proof: $L$ is
linear, so $L(y_p+y_h)=\vec b+\vec 0$. Once you see that $\tfrac{d^2}{dt^2}+1$ is a linear
operator with a $2$-dimensional nullspace, the ODE recipe and the matrix recipe are
literally the same recipe. This is the bridge to Chapter~25.
\end{connectionbox}

\begin{applicationbox}
\textbf{Underdetermined systems in machine learning and imaging.} Modern problems are
routinely \emph{wide}: you have far more unknowns $n$ than measurements $m$, so
$\mat{A}\vec{x}=\vec{b}$ has $r=m<n$ and an entire $(n-m)$-dimensional sheet of exact
solutions (Case~3). Examples: reconstructing a high-resolution image from a few sensor
readings (compressed sensing), or fitting an over-parameterized neural network where there
are more weights than training equations. Theorem~\ref{ch13:structure} says all these
solutions are equally valid \emph{as far as the equations are concerned} -- so you must add
a preference to pick one. Choosing the solution of smallest length $\norm{\vec{x}}$ gives
the \emph{minimum-norm} solution and lands you back on the row space; choosing the sparsest
$\vec{x}$ (fewest nonzeros) is the $\ell_1$ idea behind compressed sensing. Either way, you
are selecting one point $\vec{x}_p$ out of the affine sheet $\vec{x}_p+\nullsp(\mat{A})$.
The nullspace is not noise to be ignored; it is the precise space of ambiguity your prior
must resolve.
\end{applicationbox}

\section{Recipe: the complete solution in five steps}

\begin{focusbox}
Given $\mat{A}\vec{x}=\vec{b}$ with $\mat{A}\in\R^{m\times n}$:
\begin{enumerate}
  \item \textbf{Augment and reduce.} Row-reduce $[\,\mat{A}\mid\vec{b}\,]$ to a
    row echelon (or reduced) form. Identify the pivot columns; let $r$ be their number.
  \item \textbf{Check consistency.} If any zero row of $\mat{A}$ has a nonzero augmented
    entry, \textbf{stop: no solution}. Otherwise continue.
  \item \textbf{Particular solution.} Set all free variables to $0$ and solve for the pivot
    variables. Call the result $\vec{x}_p$.
  \item \textbf{Special solutions.} For each free variable, set it to $1$ and the other
    free variables to $0$; solve $\mat{A}\vec{x}=\vec{0}$ for the pivot variables. This
    gives $\vec{s}_1,\dots,\vec{s}_{n-r}$, a basis of the nullspace.
  \item \textbf{Assemble.} The complete solution is
    $\vec{x}=\vec{x}_p+c_1\vec{s}_1+\cdots+c_{n-r}\vec{s}_{n-r}$, $c_i\in\R$.
\end{enumerate}
If $r=n$ there are no free variables and $\vec{x}=\vec{x}_p$ is unique. If $r=m$ the
consistency check in step~2 can never fail.
\end{focusbox}

\begin{example}[A clean square singular system, start to finish]
\label{ch13:square}
\[
  \mat{A}=\begin{bmatrix}1&2&1\\2&4&3\\1&2&2\end{bmatrix},\qquad
  \vec{b}=\colvec{2\\5\\3}.
\]
Reduce $[\,\mat A\mid\vec b\,]$. Subtract $2R_1$ from $R_2$ and $R_1$ from $R_3$:
\[
  \left[\begin{array}{ccc|c}1&2&1&2\\2&4&3&5\\1&2&2&3\end{array}\right]
  \longrightarrow
  \left[\begin{array}{ccc|c}1&2&1&2\\0&0&1&1\\0&0&1&1\end{array}\right]
  \xrightarrow{R_3-R_2}
  \left[\begin{array}{ccc|c}1&2&1&2\\0&0&1&1\\0&0&0&0\end{array}\right].
\]
Pivots in columns $1$ and $3$; column $2$ is free, so $r=2$, and the last row reads $0=0$
-- consistent. \emph{Particular} (set $x_2=0$): from $R_2$, $x_3=1$; from $R_1$,
$x_1+2(0)+1=2\Rightarrow x_1=1$. So $\vec{x}_p=\colvec{1\\0\\1}$. \emph{Special} (set
$x_2=1$, $\vec b\to\vec 0$): $x_3=0$ and $x_1+2(1)+0=0\Rightarrow x_1=-2$, giving
$\vec{s}=\colvec{-2\\1\\0}$. Complete solution:
\[
  \vec{x}=\colvec{1\\0\\1}+c\colvec{-2\\1\\0},\qquad c\in\R.
\]
A line of solutions in $\R^3$ -- this square matrix had $r=2<3$, dependent columns, so it
is squarely in Case~4 despite being $3\times 3$.
\end{example}

\section{Counting, sanity checks, and what can go wrong}

A few quick mental checks save you from most errors.

\begin{itemize}
  \item \textbf{Dimension check.} The number of special solutions must equal the number of
    free columns, $n-r$. If you produced three special solutions but only two columns were
    free, you double-counted.
  \item \textbf{Plug $\vec{x}_p$ back in.} It costs one matrix-vector product and catches
    arithmetic slips immediately. $\mat{A}\vec{x}_p$ must equal $\vec{b}$ on the nose.
  \item \textbf{Plug each $\vec{s}_i$ in.} Each must give $\mat{A}\vec{s}_i=\vec{0}$. If
    not, you mishandled a sign during back-substitution.
  \item \textbf{Consistency is about $\vec{b}$ only.} Changing $\vec{b}$ never changes the
    nullspace or the special solutions -- those come from $\mat{A}$ alone. So if you solve
    $\mat{A}\vec{x}=\vec{b}$ for several right-hand sides, you reduce $\mat A$ once and
    reuse the special solutions every time; only $\vec{x}_p$ and the consistency verdict
    change.
\end{itemize}

\begin{walkbox}
Here is a puzzle to chew on. You are told $\mat{A}$ is $4\times6$ and that
$\mat{A}\vec{x}=\vec{b}$ has a solution for the specific $\vec{b}=\colvec{1\\1\\1\\1}$, and
that the solution is \emph{not} unique. What is the smallest the solution set's dimension
could be? Could the solution ever be unique for some \emph{other} right-hand side while
staying non-unique for this one? (Think about whether ``unique'' is a property of
$\mat{A}$ or of $\vec{b}$.)
\end{walkbox}

\begin{aestheticbox}
Step back and admire the division of labor. Two subspaces attached to $\mat{A}$ -- the
column space living in the output $\R^m$ and the nullspace living in the input $\R^n$ --
between them answer \emph{every} question you can ask about $\mat{A}\vec{x}=\vec{b}$. The
column space decides \emph{whether}; the nullspace decides \emph{how many}; and the single
number $r$ ties them together through $r=\dim\col$ and $n-r=\dim\nullsp$. When you meet the
four fundamental subspaces and the rank theorem in the next chapters, you will see this same
$r$ orchestrating the row space and left nullspace too. The complete solution is your first
full sighting of that hidden organizing number.
\end{aestheticbox}

\begin{historybox}
The decomposition of an inhomogeneous problem into a particular solution plus all solutions to the homogeneous equation is one of the grand unifying themes of mathematical analysis. It was first discovered in the eighteenth century by Leonhard Euler, Joseph-Louis Lagrange, and Jean le Rond d'Alembert while solving linear differential equations of the form $L(y) = f(t)$. They noticed that if $y_p$ is any single function satisfying $L(y_p) = f(t)$, and if $y_h$ satisfies the zero-forcing equation $L(y_h) = 0$, then linearity guarantees that $L(y_p + y_h) = L(y_p) + L(y_h) = f(t) + 0 = f(t)$. 

When nineteenth-century mathematicians began studying systems of linear algebraic equations with rectangular matrices, they found the exact same structural law: $\vec{x} = \vec{x}_p + \vec{x}_n$. The solution set of $\mat{A}\vec{x}=\vec{b}$ is not a subspace (it fails closure: adding two solutions produces $\mat{A}(\vec{x}_1+\vec{x}_2) = 2\vec{b} \ne \vec{b}$, and $\vec{0}$ is not in the set unless $\vec{b}=\vec{0}$). Instead, it is an \emph{affine flat}---a copy of the nullspace $\nullsp(\mat{A})$ displaced away from the origin by the particular anchor $\vec{x}_p$.

For decades, early algebra courses treated $m\times n$ systems with confusing lists of case-by-case rules: ``equations outnumbering unknowns'' versus ``unknowns outnumbering equations.'' Students and engineers routinely fell into the naive trap of assuming that fewer equations than unknowns guaranteed a solution, or that equal numbers guaranteed uniqueness. The modern geometric synthesis swept away this folklore. Solvability does not depend on the superficial count of equations, but strictly on whether the right-hand side lives in the column space $\col(\mat{A})$. Moreover, beginners often worry that their choice of $\vec{x}_p$ is ``wrong'' if it differs from a textbook answer. But setting free variables to zero is merely an arbitrary, convenient anchor; any vector on the affine hyperplane serves equally well as $\vec{x}_p$.
\end{historybox}

\section{Summary}

\begin{itemize}
  \item Every solution of $\mat{A}\vec{x}=\vec{b}$ is one particular solution plus a
    nullspace vector: $\vec{x}=\vec{x}_p+\vec{x}_n$. The solution set is the nullspace
    translated to pass through $\vec{x}_p$ -- an affine subspace of dimension $n-r$.
  \item \textbf{Existence} $\iff \vec{b}\in\col(\mat{A})$; detectable as ``no zero-row
    contradiction'' after augmented elimination. Existence for \emph{all} $\vec{b}$ means
    full row rank $r=m$.
  \item \textbf{Uniqueness} $\iff \nullsp(\mat{A})=\{\vec{0}\}$, i.e.\ no free variables,
    i.e.\ full column rank $r=n$.
  \item The number of solutions is always $0$, $1$, or $\infty$ -- never any other finite
    count.
  \item The four cases sort by $(r=m?)$ and $(r=n?)$: square-invertible (exactly one for
    all $\vec b$), tall full-column-rank ($0$ or $1$), wide full-row-rank ($\infty$ for
    all $\vec b$), and the general case ($0$ or $\infty$).
  \item The five-step recipe -- reduce, check consistency, particular, special, assemble --
    produces the complete solution from a \emph{single} elimination.
\end{itemize}

\begin{connectionbox}
We now have a complete constructive algorithm to describe the entire infinite solution set of any consistent rectangular system $\mat{A}\vec{x}=\vec{b}$. From a single elimination run, we identify consistency, anchor at a particular solution $\vec{x}_p$, and span every degree of freedom with $n-r$ special solutions along $\nullsp(\mat{A})$.

Yet a natural question remains: while our algorithm produces $n-r$ special solutions, what mathematically guarantees that this count is minimal? Could a clever alternative elimination path span the same nullspace with fewer vectors? What does it rigorously mean for a set of vectors to have no redundancies? Chapter~14 answers this foundational question by establishing the rigorous theory of linear independence, basis, and dimension, proving via the Steinitz exchange lemma that the number of degrees of freedom is an absolute geometric invariant of the subspace itself.
\end{connectionbox}

\section{Exercises}

\begin{exercise}[Warmup: read the verdict]
\label{ch13:ex-verdict}
For each reduced augmented matrix, state whether the system is consistent, and if so how
many solutions it has.
\[
\text{(a)}\;\left[\begin{array}{cc|c}1&2&3\\0&0&0\end{array}\right]\quad
\text{(b)}\;\left[\begin{array}{cc|c}1&0&5\\0&1&2\end{array}\right]\quad
\text{(c)}\;\left[\begin{array}{ccc|c}1&0&4&1\\0&1&2&3\\0&0&0&7\end{array}\right]
\]
\end{exercise}

\begin{exercise}[Particular plus special]
\label{ch13:ex-pns}
Find the complete solution of
$\begin{bmatrix}1&1&2\\2&2&5\end{bmatrix}\vec{x}=\colvec{3\\7}$. Identify $\vec{x}_p$, the
special solution(s), and the dimension of the solution set.
\end{exercise}

\begin{exercise}[Consistency depends on $\vec b$]
\label{ch13:ex-bvary}
Let $\mat{A}=\begin{bmatrix}1&2\\2&4\\-1&-2\end{bmatrix}$. Describe \emph{all} vectors
$\vec{b}\in\R^3$ for which $\mat{A}\vec{x}=\vec{b}$ is consistent. (Hint: what is
$\col(\mat A)$, and what linear conditions on $b_1,b_2,b_3$ cut it out?)
\end{exercise}

\begin{exercise}[A square singular system]
\label{ch13:ex-sq}
For $\mat{A}=\begin{bmatrix}1&1&1\\1&2&3\\2&3&4\end{bmatrix}$, decide which case
(of the four) it is, then give the complete solution of $\mat{A}\vec{x}=\colvec{0\\1\\1}$.
\end{exercise}

\begin{exercise}[Full column rank]
\label{ch13:ex-fcr}
The matrix $\mat{A}=\begin{bmatrix}1&1\\1&2\\1&3\end{bmatrix}$ has full column rank. For
which of $\vec{b}=\colvec{2\\3\\4}$ and $\vec{b}=\colvec{2\\3\\5}$ is the system
consistent? When it is, find the unique solution. When it is not, explain geometrically
what fails.
\end{exercise}

\begin{exercise}[Full row rank, every $\vec b$]
\label{ch13:ex-frr}
Show that $\mat{A}=\begin{bmatrix}1&0&2\\0&1&-1\end{bmatrix}$ has a solution for every
$\vec{b}\in\R^2$, and write the complete solution of $\mat{A}\vec{x}=\colvec{1\\1}$ as a
particular solution plus a line.
\end{exercise}

\begin{exercise}[Why only $0$, $1$, or $\infty$]
\label{ch13:ex-trichotomy}
Prove that a system $\mat{A}\vec{x}=\vec{b}$ cannot have exactly two distinct solutions.
(Hint: if $\vec{x}_1\neq\vec{x}_2$ both solve it, consider
$\vec{x}_1+t(\vec{x}_2-\vec{x}_1)$ for $t\in\R$ and use linearity.)
\end{exercise}

\begin{exercise}[Reusing one elimination]
\label{ch13:ex-reuse}
Let $\mat{A}=\begin{bmatrix}1&3\\2&6\end{bmatrix}$. Find the nullspace once. Then, without
re-reducing $\mat{A}$, determine the complete solution set for
$\vec{b}=\colvec{1\\2}$ and for $\vec{b}=\colvec{1\\3}$. Explain why the special solution is
identical in both cases.
\end{exercise}

\begin{exercise}[The shape from rank alone]
\label{ch13:ex-shape}
A matrix $\mat{A}$ is $5\times7$ with rank $r=4$. (a) What is the dimension of the nullspace?
(b) For a consistent $\vec{b}$, what is the dimension of the solution set? (c) Is the system
solvable for \emph{every} $\vec{b}\in\R^5$? (d) Can any right-hand side give a unique
solution? Justify each answer from the four-case framework.
\end{exercise}

\begin{exercise}[Challenge: prescribe the solution set]
\label{ch13:ex-prescribe}
Construct a $3\times3$ matrix $\mat{A}$ and a vector $\vec{b}$ so that the complete solution
of $\mat{A}\vec{x}=\vec{b}$ is exactly the line
$\colvec{1\\0\\-1}+c\colvec{2\\1\\1}$, $c\in\R$. (Hint: you need
$\colvec{2\\1\\1}$ in $\nullsp(\mat A)$ and $\mat A\colvec{1\\0\\-1}=\vec b$. Build the rows
so each is orthogonal to the nullspace direction, then read off $\vec b$.)
\end{exercise}

\begin{exercise}[Challenge: a uniqueness criterion without elimination]
\label{ch13:ex-criterion}
Suppose $\mat{A}$ is $m\times n$ and you are told that $\mat{A}\T\mat{A}$ is invertible.
Prove that $\mat{A}\vec{x}=\vec{b}$ has \emph{at most} one solution for every $\vec{b}$.
(Hint: show $\mat{A}\vec{x}=\vec 0$ forces $\vec{x}=\vec 0$ by hitting it with $\mat{A}\T$,
or by looking at $\norm{\mat A\vec x}^2=\vec x\T\mat A\T\mat A\vec x$.) Which of the four
cases does this rule out?
\end{exercise}
